% Part:sets-functions-relations
% Chapter: size-of-sets
% Section: enumerations

\documentclass[../../../include/open-logic-section]{subfiles}

\begin{document}

\olfileid{sfr}{siz}{enm}

\olsection{شمېرنې او \usetoken{S}{enumerable} سټونه}

\begin{editorial}
  دا برخه د سټونو شمېرنې څېړي او له $\PosInt$ څخه د پر هدف سټ
  بشپړو تابعو په توګه ئې تعريفوي۔ د هغو لوستونکو دپاره چې
  لږ رياضيکي شاليد لري، خبرې ګام په ګام وړاندې کوي۔ يوه
  بديله، لنډه بڼه په \olref[enm-alt]{sec} کښې ورکړل شوې، چې
  شمېرنې په بېل ډول تعريفوي: له $\Nat$ (يا د هغې له پيلنۍ
  برخې) سره د دوه اړخيزو يو پر يو تابعو په توګه۔
\end{editorial}

\begin{explain}
مخکښې مو د سټونو بېلګې د هغو د \psOblique{!!{element}s} په
ليکلو ورکړې دي۔ راځئ په زياته عمومي توګه بحث وکړو چې د يوه
سټ !!{element}s څنګه او کله په لړ کښې راوړلے شو، که هغه سټ
نامتناهي هم وي۔
\end{explain}

\begin{defn}[شمېرنه، په غير صوري ډول]
په غير صوري ډول، د سټ~$A$ \emph{شمېرنه} د~$A$ د
\psOblique{!!{element}s} يو (کېدے شي نامتناهي) لړ دے، داسې چې د $A$
هر !!{element} د لړ په يوه متناهي مقام کښې راشي۔ که $A$ شمېرنه
ولري، نو $A$ ته \emph{!!{enumerable}} ويل کېږي۔
\end{defn}

\begin{explain}
د شمېرنو په اړه څو ټکي:
\begin{enumerate}
\item يوازې هغه لړونه شمېرنې ګڼو چې پيل ولري او له لومړني
  غړي پرته د هر \psOblique{!!{element}} سمدستي
  مخکښې يو يوازينے !!{element} وي۔ په بله وينا، د لړ د لومړي
  \psOblique{!!{element}} او هر بل \psOblique{!!{element}} ترمنځ
  يوازې متناهي شمېر غړي وي۔ په ځانګړي ډول، دا معنٰی لري چې
  د شمېرنې هر !!{element} متناهي مقام لري: لومړے !!{element}
  په مقام~$1$ کښې وي، دويم په مقام~$2$ کښې، او داسې نور۔
\item د هماغه سټ~$A$ بېلې شمېرنې لرلے شو چې د
  \psOblique{!!{element}s} د راتلو په ترتيب کښې سره بېلې وي:
  $4$، $1$، $25$، $16$،~$9$ د لومړيو پنځو مربع عددونو (د سټ)
  هومره ښه شمېرنه ده لکه $1$، $4$، $9$، $16$،~$25$۔
\item تکراري شمېرنې هم شمېرنې دي: $1$، $1$، $2$، $2$، $3$،
  $3$،~\dots{} د هماغه سټ شمېرنه ده چې $1$، $2$، $3$،~\dots{}
  ئې شمېرنه کوي۔
\item کله چې يوه شمېرنه ټاکو، ترتيب او تکرار \emph{اهميت لري}:
  د مثبتو صحيح عددونو شمېرنه په $1$، $2$، $3$، $1$، \dots{}
  پيلولے شو، خو نمونه هغه وخت په اسانه ښکاري چې په معمولي
  ډول ئې په $1$، $2$، $3$، $4$،~\dots وشمېرو۔
\item شمېرنې بايد پيل ولري: \dots، $3$، $2$، $1$ د مثبتو صحيح
  عددونو شمېرنه نۀ ده، ځکه چې لومړے !!{element} نۀ لري۔ د دې
  د ليدلو دپاره چې دا له غير صوري تعريفه څنګه راځي، له ځانه
  وپوښتئ: «عدد 76 د لړ په کوم مقام کښې راځي؟»
\item لاندې لړ د مثبتو صحيح عددونو شمېرنه نۀ ده:
  $1$، $3$، $5$، \dots، $2$، $4$، $6$، \dots\@ ستونزه دا ده چې
  جفت عددونه د متناهي مقامونو پر ځاے په $\infty + 1$،
  $\infty + 2$، $\infty + 3$ مقامونو کښې راځي۔
\item تش سټ د شمېر وړ دے: تش لړ ئې شمېري!{}
\end{enumerate}
\end{explain}

\begin{prop}
  که $A$ شمېرنه ولري، بې له تکراره شمېرنه هم لري۔
\end{prop}

\begin{proof}
  فرض کړئ $A$ شمېرنه $x_1$، $x_2$، \dots{} لري، چې هر $x_i$
  د~$A$ يو !!{element} دے۔ له شمېرنې څخه د تکراري
  \psOblique{!!{element}s} په لرې کولو تکرارونه لرې کولے شو۔ د
  بېلګې په توګه، شمېرنه داسې نوې شمېرنې ته اړولے شو چې $x_i$
  پکښې هله وليکو چې د~$A$ داسې !!a{element} وي چې په $x_1$،
  \dots، $x_{i-1}$ کښې نۀ وي، يا $x_i$ له لړه لرې کړو که له
  مخکښې په $x_1$، \dots،~$x_{i-1}$ کښې راغلے وي۔
\end{proof}

وروستے استدلال ښيي چې د شمېرنو او !!{enumerable} سټونو د ښه
پوهېدو او د هغو په اړه د خبرو د ثابتولو دپاره يو زيات کره
تعريف ته اړتيا لرو۔ لاندې تعريف دا اړتيا پوره کوي۔

\begin{defn}[شمېرنه، په صوري ډول]
د سټ $A \neq \emptyset$ \emph{شمېرنه} هره !!{surjective}
تابع $f \colon \PosInt \to A$ ده۔
\end{defn}

\begin{explain}
راځئ ځان ډاډمن کړو چې صوري تعريف او هغه غير صوري تعريف چې
يو ممکن نامتناهي لړ کاروي، همارزښته دي۔ لومړے، له $\PosInt$
څخه سټ~$A$ ته هره !!{surjective} تابع د~$A$ شمېرنه کوي۔ داسې
تابع د پاسني غير صوري تعريف له مخې يوه شمېرنه ټاکي: لړ
$f(1)$، $f(2)$، $f(3)$، \dots۔ ځکه چې $f$ !!{surjective} ده،
دا تضمين شوے چې د~$A$ هر !!{element} به د کوم~$n \in \PosInt$
دپاره د~$f(n)$ قيمت وي۔ نو د $A$ هر !!{element} په لړ کښې
په کوم متناهي مقام راځي۔ ځکه چې کېدے شي تابع !!{injective}
نۀ وي، لړ تکرار لرلے شي؛ خو دا د منلو وړ دي، لکه چې پورته
ياد شول۔

بل خوا، که داسې لړ راکړل شوے وي چې د~$A$ ټول !!{element}s
شمېري، يوه !!a{surjective} تابع $f\colon \PosInt \to A$ داسې
تعريفولے شو چې $f(n)$ د لړ $n$م !!{element} وټاکو، يا که
$n$م !!{element} نۀ وي، نو د لړ وروستے !!{element} وټاکو۔
يوازې په هغه حالت کښې له دې څخه !!a{surjective} تابع نۀ
جوړېږي چې $A$ تش وي او له دې کبله لړ هم تش وي۔ نو هر ناتش
لړ يوه !!a{surjective} تابع $f\colon \PosInt \to A$ ټاکي۔
\end{explain}

\begin{defn}
  \ollabel{defn:enumerable}
  سټ~$A$ هله او يوازې هله !!{enumerable} دے چې تش وي يا شمېرنه ولري۔
\end{defn}

\begin{ex}
د مثبتو صحيح عددونو ($\PosInt$) شمېرونکې تابع يوازې د عينيت
هغه تابع ده چې په $f(n) = n$ ورکړل شوې ده۔ د طبيعي عددونو
$\Nat$ شمېرونکې تابع $g(n) = n - 1$ ده۔
\end{ex}

\begin{ex}
تابعې $f\colon \PosInt \to \PosInt$ او $g \colon \PosInt \to
\PosInt$ چې داسې ورکړل شوې دي:
\begin{align*}
f(n) & = 2n \text{ او}\\
g(n) & = 2n - 1
\end{align*}
په ترتيب سره د جفتو مثبتو صحيح عددونو او د طاقو مثبتو صحيح
عددونو شمېرنه کوي۔ خو له دوو هېڅ يوه د $\PosInt$ شمېرنه نۀ
ده، ځکه چې هېڅ يوه !!{surjective} نۀ ده۔
\end{ex}

\begin{prob}
د مثبتو مربع عددونو $1$، $4$، $9$، $16$، \dots يوه شمېرنه تعريف کړئ۔
\end{prob}

\begin{ex}
تابع $f(n) = (-1)^{n} \lceil \frac{(n-1)}{2}\rceil$ (چې
$\lceil x \rceil$ د \emph{پاسني صحيح عدد} تابع ښيي، يعنې $x$
پورته تر ټولو نږدې صحيح عدد ته ګردوي) د صحيح عددونو~$\Int$
د سټ شمېرنه کوي۔ پام وکړئ چې $f$ د مثبتو او منفي صحيح عددونو
ترمنځ په تګ راتګ سره د $\Int$ قيمتونه څنګه راوباسي:
\[
\begin{array}{c c c c c c c c}
f(1) & f(2) & f(3) & f(4) & f(5) & f(6) & f(7) & \dots \\ \\
- \lceil \tfrac{0}{2} \rceil & \lceil \tfrac{1}{2}\rceil & - \lceil \tfrac{2}{2} \rceil & \lceil \tfrac{3}{2} \rceil & - \lceil \tfrac{4}{2} \rceil  & \lceil \tfrac{5}{2}
\rceil & - \lceil \tfrac{6}{2} \rceil & \dots \\ \\
0 & 1 & -1 & 2 & -2 & 3 & -3 & \dots
\end{array}
\]
\textbf{د ژباړې د نسخې سمون (OLSIZ-001):} د انګرېزي جدول سرليک
او د فارمول قطار f(7) لري، خو د قيمتونو قطار ئې اړين قيمت -3
پرېښے و؛ دلته د هماغه فارمول له مخې ورزيات شوے دے۔
$f$ د حالتونو له مخې داسې تعريف شوې هم ګڼلے شئ:
\[
f(n) = \begin{cases}
  0 & \text{که $n = 1$ وي}\\
  n/2 & \text{که $n$ جفت وي}\\
  -(n-1)/2 & \text{که $n$ طاق او $>1$ وي}
  \end{cases}
\]
\end{ex}

\begin{prob}
  وښيئ چې که $A$ او $B$ !!{enumerable} وي، نو $A \cup B$ هم
  داسې دے۔ د دې دپاره فرض کړئ چې !!{surjective} تابعې
  $f\colon \PosInt \to A$ او $g\colon \PosInt \to B$ شته؛
  يوه !!a{surjective} تابع~$h\colon \PosInt \to A \cup B$
  تعريف کړئ او ثابته کړئ چې !!{surjective} ده۔ هغه حالتونه
  هم په پام کښې ونيسئ چې $A$ يا~$B = \emptyset$ وي۔
\end{prob}
  
\begin{prob}
  وښيئ چې که $B \subseteq A$ او $A$ !!{enumerable} وي، نو~$B$
  هم داسې دے۔ د دې دپاره فرض کړئ چې يوه !!a{surjective}
  تابع $f\colon \PosInt \to A$ شته۔ يوه !!a{surjective}
  تابع~$g\colon \PosInt \to B$ تعريف کړئ او ثابته کړئ چې
  !!{surjective} ده۔ که $B = \emptyset$ وي، څه پېښېږي؟
\end{prob}
    
\begin{prob}
  پر $n$ د استقرا په وسيله وښيئ چې که $A_1$، $A_2$، \dots،
  $A_n$ ټول !!{enumerable} وي، نو $A_1 \cup \dots \cup A_n$
  هم داسې دے۔ دا حقيقت فرضولے شئ چې که دوه سټونه $A$ او~$B$
  !!{enumerable} وي، نو~$A \cup B$ هم داسې دے۔
\end{prob}

که څه هم د يوه سټ د \psOblique{!!{element}s} د لړ کولو پر وخت
ښايي له $1$م \psOblique{!!{element}} څخه شمېر پيلول زيات طبيعي
ښکاري، رياضي پوهان د څيزونو د شمېرلو دپاره طبيعي عددونه~$\Nat$
کارول خوښوي۔ دوے د لړ د $0$م، $1$م، $2$م، او داسې نورو
\psOblique{!!{element}s} خبرې کوي۔ له دې سره سم، شمېرنه له
$\Nat$ څخه~$A$ ته د يوې \psOblique{!!a{surjective}} تابعې په توګه تعريفولے
شو۔ البته، دواړه تعريفونه همارزښته دي۔

\begin{prop}\ollabel{prop:enum-shift}
  يوه !!a{surjection} $f\colon \PosInt \to A$ هله او يوازې هله
  شته چې يوه !!a{surjection} $g\colon \Nat \to A$ شته وي۔
\end{prop}

\begin{proof}
  د ورکړل شوې \psOblique{!!a{surjection}} $f\colon \PosInt \to A$ دپاره،
  د ټولو $n \in \Nat$ دپاره $g(n) = f(n+1)$ تعريفولے شو۔ په
  اسانه ښکاري چې $g\colon \Nat \to A$ يوه !!{surjective} تابع
  ده۔ برعکس، د ورکړل شوې \psOblique{!!a{surjection}} $g\colon \Nat \to A$
  دپاره، $f(n) = g(n-1)$ تعريف کړئ۔
\end{proof}

له دې لاندې پايله ترلاسه کېږي:

\begin{cor}\ollabel{cor:enum-nat}
سټ $A$ هله او يوازې هله !!{enumerable} دے چې تش وي يا يوه
!!a{surjective} تابع $f\colon \Nat \to A$ شته وي۔
\end{cor}

پورته مو بحث وکړ چې د سټ~$A$ د \psOblique{!!{element}s} لړ په
بې تکراره لړ اړول کېدے شي۔ دا د شمېرنو دپاره هم رښتيا ده، خو
په کره ډول ئې بيانول او ثابتول لږ ګران دي۔ هره تابع
$f\colon \PosInt \to A$ بايد د ټولو $n \in \PosInt$ دپاره
تعريف شوې وي۔ که په~$A$ کښې يوازې متناهي شمېر !!{element}s وي،
نو په ښکاره داسې تابع نۀ شو لرلے چې د~$\PosInt$ په نامتناهي
شمېر \psOblique{!!{element}s} تعريف شوې وي، د~$A$ ټول
!!{element}s د قيمتونو په توګه واخلي، خو هېڅکله يو قيمت دوه
ځله وانخلي۔ په دې حالت کښې، يعنې کله چې بې تکراره لړ متناهي
وي، بايد د~$f$ دپاره د تعريف يوه بېله ساحه وټاکو، چې يوازې
متناهي شمېر !!{element}s ولري۔ د تکرارونو نۀ لرل دا معنٰی
لري چې $f$ بايد !!{injective} وي۔ ځکه چې !!{surjective} هم ده،
نو د کوم متناهي سټ $\{1, \dots, n\}$ يا $\PosInt$ او~$A$
ترمنځ يوه !!a{bijection} لټوو۔

\begin{prop}\ollabel{prop:enum-bij}
که $f\colon \PosInt \to A$ يوه !!{surjective} تابع وي (يعنې
د~$A$ شمېرنه وي)، نو يوه !!a{bijection} $g\colon Z \to A$
شته چې $Z$ يا~$\PosInt$ وي يا د کوم~$n \in \PosInt$ دپاره
$\{1, \dots, n\}$ وي۔
\end{prop}

\begin{proof}
  تابع $g$ په بازګشتي ډول تعريفوو: فرض کړئ $g(1) = f(1)$۔ که
  $g(i)$ له مخکښې تعريف شوے وي، نو $g(i+1)$ د $f(1)$،
  $f(2)$، \dots{} لومړے هغه قيمت وټاکئ چې تر اوسه په $g(1)$،
  \dots، $g(i)$ کښې نۀ وي، که داسې قيمت شته وي۔ که $A$ يوازې
  $n$ !!{element}s ولري، نو $g(1)$، \dots، $g(n)$ ټول تعريف
  شوي دي، او په دې ډول مو يوه تابع $g\colon \{1, \dots, n\}
  \to A$ تعريف کړې ده۔ که $A$ نامتناهي شمېر !!{element}s
  ولري، نو د هر $i$ دپاره بايد د $f(1)$، $f(2)$، \dots په
  شمېرنه کښې د~$A$ يو داسې !!a{element} وي چې تر اوسه په
  $g(1)$، \dots، $g(i)$ کښې نۀ وي۔ په دې حالت کښې مو يوه
  تابع $g\colon \PosInt \to A$ تعريف کړې ده۔
  
  تابع $g$ !!{surjective} ده، ځکه چې د~$A$ هر غړے په $f(1)$،
  $f(2)$، \dots{} کښې دے (ځکه چې $f$ !!{surjective} ده) او له
  دې کبله به بالاخره د کوم~$i$ دپاره د~$g(i)$ قيمت شي۔ دا
  !!{injective} هم ده، ځکه چې که داسې $j < i$ واے چې
  $g(j) = g(i)$ واے، نو $g(i)$ به له مخکښې په $g(1)$،
  \dots، $g(i-1)$ کښې و، چې زموږ د~$g$ له تعريف سره په ټکر کښې دے۔
\end{proof}

\begin{cor}\ollabel{cor:enum-nat-bij}
سټ $A$ هله او يوازې هله !!{enumerable} دے چې تش وي يا يوه
!!a{bijection} $f\colon N \to A$ شته وي، چې يا $N = \Nat$
وي يا د کوم $n \in \Nat$ دپاره $N = \{0, \dots, n\}$ وي۔
\end{cor}

\begin{proof}
$A$ هله او يوازې هله !!{enumerable} دے چې $A$ تش وي يا يوه
!!a{surjective} $f\colon \PosInt \to A$ شته وي۔ د
\olref{prop:enum-bij} له مخې، وروستۍ خبره هله او يوازې هله
سمه ده چې يوه !!a{bijective} تابع~$f\colon Z \to A$ شته وي،
چې $Z = \PosInt$ يا د کوم $n \in \PosInt$ دپاره
$Z = \{1, \dots, n\}$ وي۔ د \olref{prop:enum-shift} د ثبوت
د هماغه استدلال له مخې، دا بيا هله او يوازې هله سمه ده چې
يوه !!a{bijection} $g\colon N \to A$ شته وي، چې يا $N = \Nat$
وي يا $N = \{0, \dots, n-1\}$ وي۔
\end{proof}

\begin{prob}
  د \olref[sfr][siz][enm]{defn:enumerable} له مخې، سټ $A$ هله
  او يوازې هله د شمېر وړ دے چې $A = \emptyset$ وي يا يوه
  !!a{surjective} $f\colon \PosInt \to A$ شته وي۔ «!!{enumerable}
  سټ» په کره ډول داسې هم تعريفولے شو: سټ هله او يوازې هله
  د شمېر وړ دے چې يوه !!a{injective} تابع $g\colon A \to \PosInt$
  شته وي۔ وښيئ چې تعريفونه همارزښته دي؛ يعنې وښيئ چې يوه
  !!a{injective} تابع $g\colon A \to \PosInt$ هله او يوازې هله
  شته چې يا $A = \emptyset$ وي يا يوه !!a{surjective}
  $f\colon \PosInt \to A$ شته وي۔
\end{prob}

\end{document}
